\chapter{Tecnologías y herramientas}
\label{chap:tecnologias}

\lettrine{E}{l} desarrollo de un sistema integrado para interfaces cerebro-computadora requiere la selección cuidadosa de tecnologías y herramientas que proporcionen las capacidades necesarias para captura, procesamiento y análisis de señales biomédicas en tiempo real. Este capítulo presenta las tecnologías empleadas en el desarrollo del sistema, justificando las decisiones técnicas tomadas y describiendo cómo contribuyen al cumplimiento de los objetivos del proyecto.

\section{Criterios de selección tecnológica}
\label{sec:criterios_seleccion}

Las decisiones tecnológicas del proyecto se fundamentan en criterios específicos que priorizan la efectividad para los objetivos del proyecto, la sostenibilidad a largo plazo, y la facilidad de adopción por parte de la comunidad científica objetivo.

El rendimiento en tiempo real constituye un criterio fundamental, ya que las aplicaciones \acrshort{bci} requieren latencias mínimas para proporcionar \textit{feedback} efectivo. La integración con el ecosistema científico existente es prioritaria para facilitar la adopción del sistema por parte de investigadores que ya utilizan herramientas establecidas. La mantenibilidad y extensibilidad del código son fundamentales para la sostenibilidad a largo plazo del proyecto.

\section{Lenguajes de programación}
\label{sec:lenguajes_programacion}

La selección de lenguajes de programación constituye una decisión fundamental que afecta tanto el rendimiento como la mantenibilidad del sistema desarrollado.

\subsection{Python como lenguaje principal}
\label{subsec:python_principal}

Python se seleccionó como lenguaje principal del proyecto debido a su prevalencia en el ecosistema de computación científica~\cite{gramfort2013meg,delorme2004eeglab}. El ecosistema científico de Python proporciona bibliotecas fundamentales como NumPy para computación matricial, SciPy para algoritmos científicos, y matplotlib para visualización~\cite{van2011numpy,virtanen2020scipy}.

La integración natural con bibliotecas especializadas en neurociencia como MNE-Python y BrainFlow constituye una ventaja decisiva~\cite{gramfort2013meg,brainflow}. El soporte moderno para tipado estático mediante \textit{type hints} permite combinar flexibilidad con verificación de tipos. Como alternativa se consideró Julia por su rendimiento superior, pero la complejidad de compilación \textit{just-in-time} y las dificultades para desarrollo iterativo rápido llevaron a la selección de Python.

\subsection{JavaScript para interfaces web}
\label{subsec:javascript_web}

JavaScript ES6+ se emplea en el desarrollo del módulo web para la exploración y pruebas del API, proporcionando capacidades de visualización interactiva y comunicación en tiempo real con el \textit{backend}. La implementación utiliza una arquitectura modular basada en \gls{managers} especializados que encapsulan funcionalidades específicas como comunicación WebSocket y visualización de gráficos.

\section{\textit{Frameworks} y bibliotecas principales}
\label{sec:frameworks_principales}

La selección de \textit{frameworks} y bibliotecas principales se basa en criterios de rendimiento, facilidad de desarrollo y compatibilidad con el ecosistema científico existente. Cada herramienta seleccionada aporta capacidades específicas que contribuyen al cumplimiento de los objetivos del proyecto.

\subsection{PyQt5 para interfaces gráficas}
\label{subsec:pyqt5_interfaces}

PyQt5~\cite{pyqt5_docs} constituye la base tecnológica para la interfaz gráfica del sistema, proporcionando \gls{widget} nativos del sistema operativo con rendimiento superior y capacidades avanzadas de visualización. PyQt5 ofrece mejor rendimiento que alternativas como Tkinter para aplicaciones que requieren actualización frecuente de elementos gráficos~\cite{gui_performance_python}.

La arquitectura de señales y \gls{slots} de Qt proporciona un mecanismo robusto para la comunicación entre componentes de la interfaz de usuario. Las capacidades se complementan con PyQtGraph para visualizaciones científicas de alto rendimiento. Como alternativa se evaluó Tkinter por su simplicidad, pero las limitaciones de rendimiento para visualizaciones en tiempo real hicieron inviable su uso, derivando en la adopción de PyQt5. No se utilizó PyQt6 debido a que no estaba tan asentado como PyQt5 en el momento del desarrollo.

\subsection{FastAPI para servicios web}
\label{subsec:fastapi_servicios}

FastAPI~\cite{fastapi_benchmark} se seleccionó como base tecnológica para el desarrollo de las API REST y servicios web del sistema debido a su combinación única de alto rendimiento y facilidad de desarrollo. La generación automática de documentación OpenAPI/Swagger elimina la necesidad de mantener documentación separada del código. La validación automática de datos mediante Pydantic proporciona garantías de tipo sin código adicional.

\subsection{Typer y Rich para interfaces de línea de comandos}
\label{subsec:typer_rich_cli}

Typer~\cite{typer_docs} y Rich~\cite{rich_docs} proporcionan la base tecnológica para las interfaces de línea de comandos del sistema. Typer aprovecha las anotaciones de tipo de Python para generar automáticamente parsers de argumentos, mientras que Rich proporciona capacidades avanzadas de formateo de salida incluyendo tablas, barras de progreso y coloreado sintáctico.

\section{Bibliotecas de procesamiento científico}
\label{sec:bibliotecas_cientifico}

\subsection{NumPy y SciPy: fundamentos del procesamiento numérico}
\label{subsec:numpy_scipy}

NumPy y SciPy constituyen los cimientos del procesamiento numérico en el sistema~\cite{van2011numpy,virtanen2020scipy}. NumPy proporciona arrays multidimensionales con operaciones vectorizadas implementadas en C, mientras que SciPy extiende las capacidades con algoritmos científicos especializados para filtros digitales, transformadas tiempo-frecuencia y análisis estadístico.

\subsection{MNE-Python para análisis neurocientífico avanzado}
\label{subsec:mne_python}

MNE-Python~\cite{gramfort2013meg} se integra como biblioteca especializada para análisis neurocientífico avanzado, proporcionando algoritmos específicamente diseñados para el procesamiento de señales EEG. La biblioteca incluye implementaciones de ICA para separación de fuentes, análisis tiempo-frecuencia mediante transformadas \textit{wavelet}, y soporte nativo para formatos estándar como BIDS y EDF+.

\section{Tecnologías de captura de entrada}
\label{sec:tecnologias_captura}

La captura precisa de eventos de entrada del usuario constituye uno de los pilares fundamentales del sistema, ya que la sincronización temporal entre señales cerebrales y acciones del usuario es crítica para aplicaciones BCI efectivas.

\subsection{pynput para captura de teclado y ratón}
\label{subsec:pynput_captura}

pynput~\cite{pynput_docs} constituye la tecnología fundamental para la captura de eventos de teclado y ratón, proporcionando una interfaz unificada y multiplataforma para el acceso a eventos de entrada de bajo nivel. La capacidad multiplataforma garantiza funcionamiento consistente en Windows, macOS y Linux, utilizando las API nativas de cada sistema operativo para acceder a eventos con mínima latencia.

\subsection{PySDL2 para controladores de juego}
\label{subsec:pysdl2_controladores}

PySDL2~\cite{pysdl2_docs} proporciona soporte avanzado para controladores de juego, específicamente implementado para el controlador PlayStation 4 DualShock. La implementación utiliza la biblioteca SDL2 subyacente, ampliamente utilizada en la industria de videojuegos, que proporciona soporte robusto para una gran variedad de dispositivos de entrada.

\section{Formatos de exportación de datos}
\label{sec:formatos_exportacion}

La capacidad de exportar datos en formatos estándar reconocidos por la comunidad científica es fundamental para garantizar la interoperabilidad del sistema con herramientas existentes y facilitar la reproducibilidad de experimentos.

\subsection{BIDS: \textit{Brain Imaging Data Structure}}
\label{subsec:bids_formato}

Brain Imaging Data Structure (BIDS)~\cite{bids} representa el estándar de facto para la organización de datos neurocientíficos, proporcionando una estructura jerárquica que facilita el intercambio y análisis reproducible de datos cerebrales. Los metadatos JSON integrados permiten documentar exhaustivamente las condiciones de adquisición y metodología experimental.

\subsection{EDF+: \textit{European Data Format Plus}}
\label{subsec:edf_formato}

European Data Format Plus (EDF+)~\cite{kemp2003european} constituye un estándar ampliamente adoptado en entornos clínicos para el almacenamiento de señales biomédicas. La especificación incluye capacidades nativas para anotaciones temporales, permitiendo documentar eventos directamente en el archivo de datos.

\subsection{HDF5: \textit{Hierarchical Data Format}}
\label{subsec:hdf5_formato}

Hierarchical Data Format 5 (HDF5)~\cite{folk2011overview} proporciona capacidades avanzadas de almacenamiento para datos científicos de gran volumen, ofreciendo compresión integrada y estructura jerárquica flexible. La estructura permite organizar datos de múltiples modalidades en un único archivo con relaciones claramente definidas.

\subsection{CSV: máxima compatibilidad}
\label{subsec:csv_formato}

Comma-Separated Values (CSV) proporciona compatibilidad universal para intercambio de datos, garantizando que los datos generados puedan ser procesados por prácticamente cualquier herramienta de análisis. La simplicidad del formato facilita la inspección manual de datos y el desarrollo de scripts de procesamiento personalizados.

\section{Tecnologías de comunicación y serialización}
\label{sec:comunicacion_serializacion}

La comunicación eficiente y la serialización optimizada de datos constituyen aspectos críticos del sistema, considerando los requisitos de tiempo real y los volúmenes significativos de datos que deben transmitirse entre componentes.

\subsection{WebSocket para comunicación en tiempo real}
\label{subsec:websocket_tiempo_real}

WebSocket~\cite{websocket_rfc} proporciona comunicación bidireccional de baja latencia entre el \textit{backend} del sistema y las interfaces web, permitiendo visualización y control en tiempo real a través de navegadores sin instalación de \textit{software} adicional.

\subsection{orjson para serialización optimizada}
\label{subsec:orjson_serializacion}

orjson~\cite{orjson_docs} proporciona serialización JSON optimizada con rendimiento superior al módulo \textit{json} estándar de Python, característica crítica para la transmisión eficiente de grandes volúmenes de datos de señales EEG. El soporte nativo para tipos de datos adicionales incluyendo \textit{datetime} y arrays NumPy elimina la necesidad de conversiones manuales.

\section{Herramientas de desarrollo y calidad de código}
\label{sec:herramientas_calidad}

El mantenimiento de la calidad del código es fundamental para la sostenibilidad y confiabilidad del sistema desarrollado. Las herramientas seleccionadas proporcionan verificación automática de estándares de código y detección temprana de problemas potenciales.

\subsection{Sistema de control de calidad integrado}
\label{subsec:control_calidad_integrado}

El sistema implementa un control de calidad de código integrado basado en \textit{pre-commit hooks} que ejecuta automáticamente verificaciones antes de cada \textit{commit}. Black~\cite{black_docs} proporciona formateo automático de código Python, mypy~\cite{mypy_docs} realiza verificación de tipos estáticos, y pydocstyle~\cite{pydocstyle_docs} verifica la conformidad de \textit{docstrings} con estándares establecidos.

\subsection{Tipado estático avanzado}
\label{subsec:tipado_estatico}

La implementación extensiva de \textit{type hints} en Python facilita el desarrollo mediante autocompletado mejorado en IDEs y verificación automática de tipos~\cite{mypy_docs}. Las anotaciones de tipo sirven como documentación que puede ser verificada automáticamente, facilitando el mantenimiento del código.

\section{Resumen de tecnologías seleccionadas}
\label{sec:resumen_tecnologias}

El conjunto de tecnologías seleccionadas proporciona las herramientas necesarias para el desarrollo del sistema, equilibrando rendimiento, facilidad de desarrollo y compatibilidad con el ecosistema científico existente. La arquitectura tecnológica resultante facilita tanto el desarrollo inicial como la evolución futura del sistema, manteniendo la flexibilidad necesaria para adaptarse a nuevos requisitos y tecnologías emergentes.

La integración de estas tecnologías crea una plataforma que aborda los objetivos específicos del proyecto mientras mantiene la extensibilidad necesaria para futuras mejoras y expansiones. La selección de cada componente tecnológico busca cumplir con los requisitos de rendimiento y funcionalidad establecidos en los objetivos del proyecto.

Para mayor detalle sobre la gestión de dependencias del sistema, consulte el Anexo~\ref{chap:dependencias_detalladas}.
